% Transcribed human main text; sources and scope in tex/007.
Let $F$ be a finite field with $q$ elements.

A set $K\subset F^n$ is called a Kakeya set if it contains a line in every direction. In other words, for every vector $a\in F^n\setminus\{0\}$, there is a vector $b$ so that the line $\{at+b\mid t\in F\}$ is contained in $K$. A trivial example of a Kakeya set is the entire vector space $F^n$. Can one find a Kakeya set significantly smaller than this?

\begin{theorem}[0.2]
A Kakeya set $K\subset F^n$ has at least $c_nq^n$ elements.
\end{theorem}
\begin{proof}
Let $K$ be a Kakeya set. If $K$ is smaller than the conclusion, let $P$ be a polynomial vanishing on $K$ of degree $<q$. Let $d$ be the degree of $P$. Write $P=P_d+Q$, where $P_d$ is the sum of monomials of degree $d$ and $Q$ is a polynomial of degree $\leq d-1$.

Let $a$ be any non-zero vector. Choose $b$ so that the line $\{at+b\mid t\in F\}$ is contained in $K$. Consider the polynomial in one variable $R(t):=P(at+b)$. The polynomial $R$ vanishes for each $t\in F$. It has degree $\leq d<q$, and so its coefficeints all vanish. In particular, its coefficient of degree $d$ vanishes. But the coefficient of $t^d$ in $R$ is exactly $P_d(a)$. So we see that $P_d(a)$ vanishes for all $a\in F^n\setminus\{0\}$. But since the degree of $P_d$ is $d<q$, it easily follows that $P_d$ vanishes at $0$ also. Then we see that $P_d$ is identically zero, and we reach a contradiction.
\end{proof}

\paragraph{所引用的维数计数结果（Lecture 1，pp.~1--2）}
\begin{corollary}[0.2]
If $\dim V(d)>s$, then there is a non-zero polynomial $P$ of degree $\leq d$ that vanishes on the given finite set.
\end{corollary}
The dimension of $V(d)$ is $\binom{d+n}{n}\geq d^n/n!$. For $n$ fixed and $d$ large, $d^n/n!$ is a good approximation of the dimension. Therefore, we get the following corollary.
\begin{corollary}[0.3]
For any set of $s$ points in $F^n$, there is a non-zero polynomial that vanishes on the set with degree $\leq ns^{1/n}$.
\end{corollary}
